\section{Illumination, Speech, and the Angelic Hierarchy}

Human beings stand aloof from one another, as St.~Gregory the Great writes in the \emph{Moralia in Iob}~(II), ``as it were behind the wall of the body,'' forced to go out ``by the door of the tongue'' whenever they wish to make themselves known. Pure spirits face no wall of flesh between mind and mind; the only veil hiding one angel's thought from another is his own free will. Turning from the angels' substance, knowledge, will, and fall (\ST{I, qq.~50--64}) to their place in divine government (\ST{I, qq.~106--114}), St.~Thomas Aquinas considers the heavenly spirits as citizens of an ordered commonwealth, united under God by downward illumination, interior speech, and hierarchical ministry---while even the fallen spirits retain a shadowed remnant of natural rank beneath divine justice.

\subsection{Angelic Illumination and the Freedom of the Will (\ST{I, q.~106})}

Because intellectual light is the manifestation of truth (Eph 5:13; 3:8), one angel enlightens another by communicating known truth downward (\ST{I, q.~106, a.~1}; \CH{VIII}). This action works through both principles of understanding. First, a superior angel strengthens the inferior angel's intellectual power by turning toward him in spiritual attention, as a hotter body intensifies the heat of a cooler one brought near it. Second, the superior divides his universal conception into more particular truths proportioned to the lower mind, as a teacher breaks a single principle into lessons for a pupil---fulfilling Pseudo-Dionysius (\CH{XV}): ``Every intellectual substance with provident power divides and multiplies the uniform knowledge bestowed on it by one nearer to God, so as to lead its inferiors upwards by analogy.'' Every blessed angel sees the divine essence immediately, so none teaches another ``Know the Lord'' (Jer 31:34); yet higher spirits behold deeper reasons (\emph{rationes}) of divine works in God and impart them downward (\ST{I, q.~106, a.~1, ad~1}; \emph{DN}~IV). No angel infuses the light of nature, grace, or glory, which flow from God alone (\ST{I, q.~106, a.~1, ad~2}), nor does angelic teaching contradict St.~Augustine's rule that God alone immediately forms the mind (\emph{Eighty-Three Diverse Questions}, q.~51), since God acts as exemplar and ultimate perfecting form while angels dispose the mind toward him (\ST{I, q.~106, a.~1, ad~3}). In \emph{De veritate} (q.~9, a.~1, ad~10), Aquinas adds that illumination infuses no new intelligible species: strengthened by superior light, the lower angel reads deeper truths within the innate species he already possesses.

Angelic power halts before the threshold of the will (\ST{I, q.~106, a.~2}). God alone, Author of intellectual nature, moves any will from within, and God alone moves it efficaciously on the side of the object by unveiling the universal Good (Ex~33:18--19). An angel inclines another's will solely from without, by intellectual persuasion. The Dionysian triad of cleansing, enlightening, and perfecting (\CH{VII}; \emph{EH}~VI) therefore belongs among angels strictly to the intellect: cleansing removes simple \emph{nescientia} (lack of knowledge), illumination imparts cognitive light, and perfection brings the mind to actual vision of truth (\ST{I, q.~106, a.~2, ad~1}; \emph{De veritate} q.~9, a.~3). Even the Seraphim kindle love in lower choirs through persuasion rather than coercion (\ST{I, q.~106, a.~2, ad~2}), while Aristotle's rule that a higher appetite moves a lower applies only to human reason commanding bodily passion (\ST{I, q.~106, a.~2, ad~3}).

Illumination flows strictly downward: an inferior angel never enlightens a superior, nor does God bypass higher spirits to illuminate lower ones independently (\ST{I, q.~106, a.~3}; \CH{IV}; \emph{EH}~V). God works miracles outside corporeal laws to draw men to faith, yet bypassing the invisible celestial hierarchy would offer no visible sign to men and serves no higher order above grace (\ST{I, q.~106, a.~3}; compare \emph{De veritate} q.~9, a.~2, quoting Dionysius's maxim \emph{omnium divinius est Dei cooperatorem fieri}, ``the most divine of all things is to become a fellow-worker with God''). In the earthly Church, sanctity and learning can diverge from ecclesiastical rank, allowing an inferior to instruct a superior (1 Cor 14:31); in heaven, grade corresponds strictly to nearness to God (\ST{I, q.~106, a.~3, ad~1}), and every angel's voluntary attention obeys the divine law establishing that order (\ST{I, q.~106, a.~3, ad~3}).

Because goodness diffuses itself (1 Pet 4:10), superior angels hold back nothing their inferiors can receive (\ST{I, q.~106, a.~4}); as St.~Gregory writes (cited through Peter Lombard, \emph{II Sent.} d.~9), ``In that heavenly country, though there are some excellent gifts, yet nothing is held individually.'' The lower spirit still receives the same truth less eminently, like a pupil learning from a master (\ST{I, q.~106, a.~4, ad~1}). Lower angels were never wholly ignorant of the Incarnation (Lombard, \emph{II Sent.} d.~11; \CH{VII}), yet they understood it less fully before its fulfillment and grew in insight as the mystery unfolded (\ST{I, q.~106, a.~4, ad~2}). Until the Day of Judgment, God continually unveils new designs of providence and salvation to the highest spirits, giving them fresh mysteries to impart down the ranks (\ST{I, q.~106, a.~4, ad~3}).

\angeldossier{\ST{I, q.~106, aa.~1--4} (Angelic illumination and the movement of the will)}{Common teaching (aa.~1--2: downward illumination and God's exclusive interior governance of the will); Thomist position (aa.~3--4: strict unidirectionality without divine bypass and proportional communication of providential knowledge).}{Eph 3:8, 5:13; Ex 33:18--19; Jer 31:34; 1 Cor 14:31; 1 Pet 4:10; Pseudo-Dionysius, \CH{IV, VII, VIII, XII, XV}, \emph{DN}~IV, \emph{EH}~V--VI; St.~Augustine, \emph{De div. quaest. LXXXIII} q.~51; St.~Gregory the Great, \emph{Hom. in Ev.}~34; Peter Lombard, \emph{II Sent.} dd.~9, 11; St.~Thomas Aquinas, \emph{De veritate} q.~9, aa.~1--3.}{In \emph{De veritate} q.~9, a.~1, ad~10, Aquinas specifies that angelic illumination infuses no new intelligible species, strengthening the inferior intellect to read deeper truths from its existing innate species.}

\subsection{The Speech of Pure Spirits (\ST{I, q.~107})}

St.~Paul's ``tongues of men and of angels'' (1 Cor 13:1) points to spiritual communication freed from vocal sound (\ST{I, q.~107, a.~1}). Following St.~Gregory (\emph{Moralia}~II), Aquinas tracks an intelligible concept across three stages: stored habitually in memory (St.~Augustine, \emph{De Trinitate} XIV.6--7), conceived actually as an ``interior word'' at the will's command, and finally directed by the will \emph{in ordine ad alium}---ordered to be made known to another. Two obstacles veil human thought: our own free will (1 Cor 2:11) and the opaque wall of our body, which requires audible signs. Having no body, an angel faces only the barrier of the will; the instant he wills his concept to be manifest to another spirit, the listener perceives it directly, so that an angel's ``tongue'' is his spiritual power of disclosing his mental word (\ST{I, q.~107, a.~1, ad~1--2}). Good angels require no signal to wake attention because they see one another continually in the Word, while even demons arrest another mind naturally by intelligible force (\ST{I, q.~107, a.~1, ad~3}).

Unlike illumination, speech moves in both directions: an inferior angel speaks upward to a superior (\ST{I, q.~107, a.~2}), as when lower choirs ask, ``Who is this King of Glory?'' (\CH{VII}). Every illumination is speech, yet not every speech is illumination. Mental concepts trace to one of two principles: God as First Truth, or the speaker's created will. Communicating truths grounded in the First Truth both speaks and illuminates, and descends strictly from above; communicating choices rooted in the created will (``I wish to learn or do this'') is speech alone, because a created will is no rule of truth, and each spirit stands supreme over his own volitions. Only in God are speech and illumination identical, since the divine will is itself the rule of truth (\ST{I, q.~107, a.~2, ad~3}; Gregory, \emph{Moralia}~II). Here the \emph{Summa} sharpens Aquinas's earlier Parisian account in \emph{De veritate} (q.~9, a.~5), which distinguished speech from illumination by object-presence versus intellectual strength, grounding the contrast instead in the causal principle of the concept.

Angels speak also to God (\ST{I, q.~107, a.~3}; Zach 1:12: ``O Lord of hosts, how long wilt thou not have mercy on Jerusalem?''). No creature informs the Omniscient, yet speech may be ordered to receiving rather than giving, like a disciple questioning a master (\ST{I, q.~107, a.~3, ad~1}). Angels address God continuously in wondering praise of his incomprehensible excellence (Gregory, \emph{Moralia}~II), and occasionally by consulting divine wisdom before new ministries (\ST{I, q.~107, a.~3, ad~2}).

Spatial distance neither delays nor impedes angelic speech (\ST{I, q.~107, a.~4}; Luke 16:24). Pure intellectual operation abstracts wholly from ``here'' and ``now,'' untied even to bodily phantasms. Speech exists where the speaking angel is---preserving St.~John of Damascus's rule that ``an angel works where he is'' (\emph{De fide orthodoxa} I.13)---yet perceiving a concept directed to oneself requires no physical nearness (\ST{I, q.~107, a.~4, ad~1}). Isaias's Seraphim ``cried one to another'' (Is 6:3) not in loud volume across space, but in grandeur of truth and ardor of love; as St.~Gregory says (\emph{Moralia}~II), ``The less one desires, the less one cries out'' (\ST{I, q.~107, a.~4, ad~2}).

Nor does independence of space make angelic speech a public broadcast (\ST{I, q.~107, a.~5}). Because speech arises from voluntary direction toward a chosen hearer, an angel can direct his mental concept to one companion while withholding it from all others (\ST{I, q.~107, a.~5, resp. and ad~1--2}). Truths descending from the First Truth are shared across the entire celestial society (\CH{XV}), whereas private speech remains bounded by the speaker's will (\ST{I, q.~107, a.~5, ad~3}).

\angeldossier{\ST{I, q.~107, aa.~1--5} (The speech of the angels)}{Common teaching (aa.~1, 3--5: interior spiritual speech, worshipful speech to God, independence of spatial distance, and volitional privacy); Thomist position (a.~2: distinction between unidirectional illumination and bidirectional speech grounded in First Truth versus created will).}{1 Cor 2:11, 13:1; Zach 1:12; Is 6:3; Luke 16:24; St.~Gregory the Great, \emph{Moralia in Iob}~II, XVIII; Pseudo-Dionysius, \CH{VII, XV}; St.~Augustine, \emph{De Trinitate} XIV.6--7; St.~John of Damascus, \emph{De fide orthodoxa} I.13; St.~Thomas Aquinas, \emph{De veritate} q.~9, aa.~4--7.}{Bl.~John Duns Scotus explains angelic speech through impressed intelligible species rather than bare volitional unveiling (a.~1), and later scholastics debated whether an unveiled concept becomes publicly visible to all angels (a.~5). In \emph{De veritate} q.~9, a.~5, Aquinas framed speech versus illumination by object-presence versus intellectual strength before sharpening the distinction causally in a.~2.}

\subsection{Hierarchies, Choirs, and the City of Angels and Saints (\ST{I, q.~108})}

Hierarchy means a ``sacred principality''---a ruler and an ordered multitude under his governance (\ST{I, q.~108, a.~1}). Viewed on the side of God the one Prince, angels and blessed men compose one commonwealth, as St.~Augustine teaches: ``There are two cities, that is, two societies, one of the good angels and men, the other of the wicked'' (\emph{City of God}~XII.1). Viewed on the side of the subjects and their modes of receiving divine light, human beings---one in species and taught under sensible veils (\CH{I})---form the ecclesiastical hierarchy, whereas angels---differentiated in species and taught in intelligible purity---form three celestial hierarchies (\CH{VI}; \ST{I, q.~108, a.~1, ad~2--3}). The first hierarchy, stationed ``in the vestibule of God'' (\CH{VII}), beholds the reasons of things directly in God as first universal principle; the second contemplates them in universal created causes; the third as applied to particular effects. Aquinas warns against calling the Blessed Trinity a ``supercelestial hierarchy,'' since the Divine Persons share a natural order of origin without hierarchical cleansing, enlightening, or perfecting (\ST{I, q.~108, a.~1}; \CH{III}).

Each angelic hierarchy contains three orders, making nine choirs in all (\ST{I, q.~108, a.~2}; Eph 1:20--21; \CH{VI}). Just as every ordered city divides into nobles, honorable middle citizens (\emph{populus honorabilis}), and common folk, each celestial hierarchy comprises summit, middle, and base. All spiritual gifts are held in common (Peter Lombard, \emph{II Sent.} d.~9), yet higher ranks possess them more eminently by communicating them to others (\ST{I, q.~108, a.~2, ad~2}). Unlike deacons, priests, and bishops in \emph{Ecclesiastical Hierarchy}~V, even the least angel excels John the Baptist (Matt 11:11) and exercises cleansing, enlightening, and perfecting in a higher mode (\ST{I, q.~108, a.~2, ad~3}).

Countless angels belong to each choir, as Isaias shows when ``the Seraphim cried to one another'' (\ST{I, q.~108, a.~3}; Is 6:3). Although each angel forms its own species (\ST{I, q.~50, a.~4}) and holds a singular office known fully to God alone, our imperfect knowledge (\CH{VI}) groups them by shared office-likeness, with first, middle, and last ranks within every choir (\ST{I, q.~108, a.~3, ad~1--2}; \CH{X}). Like points along the border of a black-and-white surface, two angels at the edge of adjacent choirs stand closer in natural proximity, yet less akin in office-fitness than companions within one order (\ST{I, q.~108, a.~3, ad~3}).

Relative to their natural end, the choirs differ by natural gifts; relative to the beatific vision, they are distinguished \emph{adequately} by gifts of grace (Rom 5:5; Augustine, \emph{City of God}~XII) and \emph{dispositively} by natural gifts, since angels receive grace in proportion to natural capacity (\ST{I, q.~108, a.~4}; \ST{I, q.~62, a.~6}; Peter Lombard, \emph{II Sent.} d.~9). Human ecclesiastical orders, by contrast, rest on gratuitous consecration rather than natural endowment.

Scripture supplies all nine names: Seraphim (Is 6:2), Cherubim (Ezech 1; 10:15--20), Thrones (Col 1:16), Dominations, Virtues, Powers, and Principalities (Eph 1:21), Archangels (Jude 9), and Angels (\ST{I, q.~108, a.~5}). In any graded series, a perfection exists \emph{by property} (proportionate to a rank), \emph{by excess} (eminently in superiors), or \emph{by participation} (partially in inferiors). Each choir takes its name from its proper excellence, while Pseudo-Dionysius (\CH{VII--IX}) emphasizes interior perfections and St.~Gregory (\emph{Homily~34 on the Gospels}, \S10) emphasizes exterior ministries. Reading both Doctors across Article~5 defines the nine choirs:

\begin{itemize}
\item \textbf{Seraphim} take their name from an excess of charity imaged by fire: perpetual upward motion toward God, penetrating heat that kindles and cleanses inferiors, and inextinguishable brightness (ad~5; \CH{VII}).
\item \textbf{Cherubim} signify ``fullness of knowledge''---clear vision of God, full reception of divine light, contemplation of divine order, and generous outpouring of wisdom (ad~5; \CH{VII}).
\item \textbf{Thrones} mirror four marks of a royal seat: elevation, divine firmness, bearing God and his judgments downward, and open readiness to receive him (ad~6; \CH{VII}).
\item \textbf{Dominations} share God's lordship (\emph{DN}~XII) through unservile freedom and the command prescribing what is to be done (ad~2--3; \CH{VIII}).
\item \textbf{Virtues}---distinct from the generic title ``heavenly virtues''---denote ``virile and immovable strength'' for fearless divine execution and ministerial miracles (ad~1; \CH{VIII}).
\item \textbf{Powers} signify sacred ordination (Rom 13:2), ordering execution and restraining hostile demons (ad~3; \CH{VIII}).
\item \textbf{Principalities} (\emph{principari}, to be first) lead sacred action over peoples and kingdoms, like precentors leading a choir (Ps 67:26 [Heb 68:25]; ad~3; \CH{IX}).
\item \textbf{Archangels} mediate between Principalities and Angels as ``angel-princes'' announcing supreme mysteries (ad~4; \CH{IX}).
\item \textbf{Angels} hold the common messenger-name as their own at the base of the ladder, announcing divine things immediately to men (ad~1; \CH{V}).
\end{itemize}

In Article~6, Aquinas confronts a classic patristic divergence. Pseudo-Dionysius (\CH{VI--IX}) and St.~Gregory (\emph{Homily~34}, \S7) agree on the first triad and the lowest two choirs, yet invert Principalities and Virtues: Dionysius sets \emph{Dominations--Virtues--Powers} and \emph{Principalities--Archangels--Angels} (matching Eph 1:20--21 ascending), whereas Gregory in \emph{Homily~34} sets \emph{Dominations--Principalities--Powers} and \emph{Virtues--Archangels--Angels} (matching Col 1:16 descending)---and in his earlier \emph{Moralia in Iob}~(XXXII.23.48) Gregory gives yet another ascending sequence placing Thrones below Dominations, Virtues, Principalities, and Powers. Keeping the promise he made in \emph{Super Ephesios}~1:21 and fulfilled in \emph{Super Colossenses}~1:16, Aquinas reconciles both schemes (\ST{I, q.~108, a.~6}). Dionysius structures the triads by the three stages of action: contemplating the End in God (Thrones receive him, Cherubim know his secrets, Seraphim unite to him in fiery love, standing closest to the Holy Spirit), disposing universal government (Dominations command, Virtues empower, Powers order), and executing particular ministries (Principalities lead nations, Dan 10:13; Archangels mediate; Angels execute). Explaining why Seraphim outrank Cherubim, \ST{I, q.~108, a.~6, ad~3} observes that for objects \emph{below} us, knowing surpasses loving, whereas for objects \emph{above} us---above all God---loving surpasses knowing because love draws the lover into God himself. Gregory's sequence follows St.~Augustine's cosmic governance (\emph{De Trinitate}~III): Dominations command, Principalities govern good spirits, Powers curb demons (Rom 13:3--4), Virtues work miracles, and Archangels and Angels announce. Thus Gregory's ``Principalities'' match Dionysius's ``Virtues,'' and Gregory's ``Virtues'' match Dionysius's ``Principalities'' (\ST{I, q.~108, a.~6, ad~4}).

These nine orders endure beyond the Day of Judgment (\ST{I, q.~108, a.~7}; Judg 5:20: ``Stars remaining in their order and courses''). Grades of nature and glory last forever; temporal offices guiding pilgrims cease when Christ hands over the kingdom (1 Cor 15:24; ad~1), while offices fitting triumph endure like military ranks in victory. Even after new revelations end, higher angels will eternally enlighten lower angels---and saved human beings---because keeping scientific knowledge of conclusions continuously depends on superior middle-term light (\ST{I, q.~108, a.~7, ad~2--3}).

Saved human beings are assumed into these nine angelic choirs according to grace and glory, never changing their human nature (\ST{I, q.~108, a.~8}; Matt 22:30; Luke 20:36). Because grace---the perfecting principle of every choir---comes to men by divine liberality rather than natural rank, redeemed souls merit degrees of glory equal to angels in every order (\ST{I, q.~108, a.~8, resp. and ad~1}). Rejecting both the denial of human equality with angels and the claim that only virgins join the nine choirs while others form a separate tenth order, Aquinas invokes St.~Augustine (\emph{City of God}~XII.9): ``there will not be two societies of men and angels, but only one; because the beatitude of all is to cleave to God alone.'' Under common law angels administer earthly affairs, whereas departed saints intervene only by special dispensation (ad~2; Augustine, \emph{De cura pro mortuis}~XIII, XVI); and St.~John Chrysostom (\emph{Homily~28 on Matthew}) refutes the pagan error that demons are souls of the dead (ad~3).

\angeldossier{\ST{I, q.~108, aa.~1--8} (The angelic degrees of hierarchies and orders)}{Common teaching (aa.~1--3, 5, 7--8: ordered angelic society under God and multiple scriptural ranks belong to Church teaching, while the Dionysian 3$\times$3 nine-choir architecture is common teaching); Thomist position (a.~4: orders distinguished adequately by grace and dispositively by nature); Disputed (a.~6: relative ranking of the middle and lower orders).}{Is 6:2--3; Ezech 1, 10; Dan 10:13; Judg 5:20; Ps 67:26 [Heb 68:25]; Matt 11:11, 22:30; Luke 20:36; Rom 5:5, 13:2--4; 1 Cor 15:24; Eph 1:20--21; Col 1:16; Heb 1:14; Jude 9; Pseudo-Dionysius, \CH{I, III, V--X}, \emph{EH}~V, \emph{DN}~XII; St.~Augustine, \emph{De civitate Dei}~XII.1, 9, \emph{De cura pro mortuis}~XIII, XVI; St.~Gregory the Great, \emph{Hom. in Ev.}~34.7, 10, \emph{Moralia}~XXXII.23.48; Peter Lombard, \emph{II Sent.} d.~9; St.~John Chrysostom, \emph{Hom. in Matt.}~28; \emph{CCC}~331.}{Pseudo-Dionysius (\CH{VI--IX}) orders the middle and lower triads as Dominations--Virtues--Powers and Principalities--Archangels--Angels; St.~Gregory the Great (\emph{Hom. in Ev.}~34.7) places Principalities above Powers and Virtues below Powers, while \emph{Moralia}~XXXII.23.48 lists Thrones beneath Dominations, Virtues, Principalities, and Powers. On a.~4, Peter Lombard (\emph{II Sent.} d.~9) stresses gifts of grace without Aquinas's dispositive proportion to distinct angelic species.}

\subsection{Rank Without Charity: The Ordering of the Fallen Angels (\ST{I, q.~109})}

Sin destroys charity without annihilating created nature. When St.~Paul warns that our struggle is ``against principalities and powers, against the rulers of the world of this darkness'' (Eph 6:12), he attests that natural rank endures among the fallen (\ST{I, q.~109, a.~1}). Demons never entered the orders of glory and fell from initial habitual grace (\ST{I, q.~62, a.~3}), yet they retain their natural orders because their natural gifts remain intact (\emph{DN}~IV). Evil subsists only in a good subject (St.~Augustine, \emph{De natura boni}~III; ad~1); on God's part their ordering is sacred, on theirs unholy (ad~2); and Scripture never names demons Seraphim, Thrones, or Dominations, since burning love, divine indwelling, and spiritual liberty contradict mortal sin (ad~3).

Because action follows nature, genuine precedence (\emph{praelatio}) obtains among demons under divine wisdom, which ``reacheth from end to end mightily, and ordereth all things sweetly'' (\ST{I, q.~109, a.~2}; Wis 8:1; Gloss on 1 Cor 15:24). Their authority rests on God's justice rather than theirs (ad~1). Though pride breeds strife (Prov 13:10), weaker demons submit to stronger ones not in friendship, but in common hatred of men and God's justice (ad~2)---and presiding in wickedness deepens only the superior demon's misery (ad~3).

Yet no illumination passes among demons (\ST{I, q.~109, a.~3}; Ecclus 34:4: ``What can be made clean by the unclean?''). True illumination manifests truth referred to God as source of all light (\ST{I, q.~107, a.~2}); since demonic malice turns away from God, demons impart concepts solely by speech (ad~1) and need no illumination for natural knowledge already complete at creation (ad~2).

Over every fallen spirit stand the holy angels (\ST{I, q.~109, a.~4}), confirming St.~Augustine (\emph{De Trinitate}~III)---``the treacherous and sinful spirit of life is ruled by the rational, pious, and just spirit of life''---and St.~Gregory (\emph{Homily~34}, \S10) on the Powers who curb hostile spirits. Authority derives from nearness to God; beatified angels communicate God's judicial decrees to demons as a judge's assessors instruct executioners---illumination on the good angel's side, darkened on the demon's (ad~1)---and permit certain harms only under divine wisdom (ad~2). Even a naturally lower holy angel presides over a naturally higher demon, because the power of divine justice to which the saintly spirit cleaves outweighs every created gift of nature (\ST{I, q.~109, a.~4, ad~3}; 1 Cor 2:15).

\angeldossiersettled{\ST{I, q.~109, aa.~1--4} (The ordering of the bad angels)}{Common teaching (natural rank and functional precedence endure among fallen spirits without holiness or beatifying illumination; good angels hold authority over demons).}{Eph 6:12; Wis 8:1; Ecclus 34:4; Prov 13:10; 1 Cor 2:15, 15:24 (Gloss); Pseudo-Dionysius, \emph{DN}~IV; St.~Augustine, \emph{De natura boni}~III, \emph{De Trinitate}~III; St.~Gregory the Great, \emph{Hom. in Ev.}~34.10.}
